\section{Adaptive architecture targeting}
\label{sec:adaptive_arch}

The original NVRTC compile sites hardcoded
\texttt{arch: Some("sm\_86")} --- the Ampere compute capability of
the development RTX A1000 / RTX 3080. This made the binary
device-specific. Running on a cluster H100 (sm\_90), an L40
(sm\_89), or a Blackwell RTX 5090 (sm\_120) required either an
explicit rebuild for that target or accepting that NVRTC would
back-target sm\_86 PTX and emit a runtime warning about
deprecated PTX features.

\subsection{Implementation}

We added \texttt{gpu\_recursive::device\_nvrtc\_arch} that queries
the device's CC at runtime via the existing
\texttt{CU\_DEVICE\_ATTRIBUTE\_COMPUTE\_CAPABILITY\_\{MAJOR,MINOR\}}
attribute pattern, builds the \texttt{sm\_\{major\}\{minor\}}
string, and passes it to both NVRTC compile sites
(\texttt{GpuRecursiveContext::new} for recursive kernels,
\texttt{GpuTransportContext::new} for flat transport).

The minimum supported compute capability is CC 6.0 / sm\_60
(Pascal). Below that the helper returns an explicit error
naming the detected capability and explaining that
\texttt{atomicAdd(double*, double)} is unavailable. Falling
through silently would compile a kernel that synthesises the
atomic via lock-cmpxchg and corrupts the spectrum-hardening
tallies under contention.

The helper is \texttt{pub} (not \texttt{pub(crate)}) so binding
crates and external benchmarks can target the detected device
without re-implementing the query. \texttt{GpuTransportContext}
caches the arch string at construction time, exposed via
\texttt{nvrtc\_arch()}, so observability paths can surface what
arch the kernels were actually compiled against.

\subsection{Cost}

One bounded $\sim$6-byte \texttt{Box::leak} per process per
context (NVRTC's \texttt{CompileOptions.arch} is
\texttt{Option<\&'static str>}; the leaked string lives for the
process lifetime, which matches when the GPU context is
released anyway). The \texttt{OnceLock}-based shared context
means this leak happens at most once per process even across
ICSBEP sweep cases.

\subsection{What this unblocks}

A cluster operator running the engine across a heterogeneous
fleet can deploy one binary; each node compiles its own PTX at
startup against its own hardware. The same applies to cloud
GPUs that come and go: an H100 spot instance does not need a
custom build, just the binary and the data.
